%%%%%%%%%%%%%%%%%%%%%%%%%%%%%%%%%%%%%%%%%
% Developer CV
% LaTeX Template
% Version 1.1 (February 24, 2025)
%
% This template originates from:
% https://www.LaTeXTemplates.com
%
% Authors:
% Jan Vorisek (jan@vorisek.me)
% Based on a template by Jan Küster (info@jankuester.com)
% Modified for LaTeX Templates by Vel (vel@LaTeXTemplates.com)
%
% License:
% The MIT License (see included LICENSE file)
%
%%%%%%%%%%%%%%%%%%%%%%%%%%%%%%%%%%%%%%%%%

%----------------------------------------------------------------------------------------
%	PACKAGES AND OTHER DOCUMENT CONFIGURATIONS
%----------------------------------------------------------------------------------------

\documentclass[9pt]{developercv} % Default font size, values from 8-12pt are recommended
\usepackage[dvipsnames]{xcolor}
\usepackage{graphicx}
\usepackage{tcolorbox}
\usepackage{enumitem}
% \setlist[itemize]{noitemsep,topsep=2pt,parsep=3pt}

\setlength{\leftmargini}{.85em} 

%----------------------------------------------------------------------------------------

\begin{document}

%----------------------------------------------------------------------------------------
%	TITLE AND CONTACT INFORMATION
%----------------------------------------------------------------------------------------

\begin{minipage}[t]{0.4\textwidth} % Left column with your name and title, change the width as needed
	\vspace{-\baselineskip} % Required for vertically aligning minipages
	
	% If your name is very short: use just one of the lines below
	% If your name is very long: reduce the font size or make the current column wider (and reduce the others proportionately)
	\colorbox{violet!70!white}{{\HUGE\textcolor{white}{\textbf{\MakeUppercase{Oleksii}}}}} % First name
	
	\colorbox{violet!90!blue}{{\HUGE\textcolor{white}{\textbf{\MakeUppercase{Silchenko}}}}} % Last name
	
	\vspace{6pt} % Vertical whitespace
	
	{\huge Software Engineer} % Career or current job title
\end{minipage}
\hfill % Automatic horizontal whitespace
\begin{minipage}[t]{0.36\textwidth} % Center column with the first column of icons
	\vspace{-\baselineskip} % Required for vertically aligning minipages
	% The first parameter is the FontAwesome icon name, the second is the box size and the third is the text
	% Other icons can be found by referring to fontawesome5.pdf (supplied with the template) and using the word after \fa in the command for the icon you want
	\icon{MapMarker}{12}{Kaunas, Lithuania}\\
	\icon{At}{12}{\href{mailto:silchenko.whatthetea@gmail.com}{silchenko.whatthetea@gmail.com}}\\	
	\icon{Linkedin}{12}{\href{https://www.linkedin.com/in/whatthetea/}{in/whatthetea}}\\
	\icon{Github}{12}{\href{https://github.com/whatthetea}{github.com/whatthetea}}\\
	\icon{Globe}{12}{\href{https://whatthetea.neocities.org/}{Personal Website (WhatTheJunkyard)}}
\end{minipage}
\hfill % Automatic horizontal whitespace
\begin{minipage}[t]{0.2\textwidth} % Right column 
	\vspace{-\baselineskip} % Required for vertically aligning minipages
	\raggedleft
	\colorbox{black}{\includegraphics[width=3cm]{profile.jpg}}
\end{minipage}
%----------------------------------------------------------------------------------------
%	INTRODUCTION
%----------------------------------------------------------------------------------------

\begin{minipage}[t]{\textwidth} % Left column with the introduction text, change the width as needed
	\vspace{-\baselineskip} 
	\cvsect{Summary}
	
	Software developer with over 6 years of programming experience, including more than 2.5 years of commercial .NET development. 
	Experienced in Windows desktop applications, Azure DevOps automation, backend services, and legacy system modernization. 
	I enjoy taking ownership of complex engineering challenges, making things simple, and sharing knowledge. 
	Those qualities stem from my passion for retro tech like 16-bit consoles or my 1987 VW Jetta Mk2 project car

	\end{minipage}
\hfill % Automatic horizontal whitespace

%----------------------------------------------------------------------------------------
%	EXPERIENCE
%----------------------------------------------------------------------------------------

\newcommand{\location}[1]{
	\small\textit{#1}
}

\cvsect{Experience}

\begin{entrylist}
	\entry
		{03.2024 -- present}
		{C\# Developer}
		{Digital Cloud Technologies}
		{\texttt{ASP.NET Core}\slashsep\texttt{Azure DevOps}\slashsep\texttt{MSBuild}
		\slashsep\texttt{Docker}\slashsep\\\texttt{Python}
		\slashsep\texttt{AI}\slashsep\texttt{PowerShell}\slashsep\texttt{WPF}\slashsep\texttt{UWP}\slashsep\texttt{WinUI}}
		{\location{Remote}}
		{
			{\bf{Web \& AI:}}

			\begin{itemize}
				\item Python FastAPI microservice for a RAG chatbot over a legal-document database: Unstructured.io, SentenceTransformers, PostgreSQL, scheduled tasks and tracked documents configurable at runtime; part of the production system
				\item Took part in maintaining distributed monolith ASP.NET Core solution with HTML jQuery views, resolved issues within AI RAG chatbot module 
			\end{itemize}
			
			{\bf{DevOps \& Developer Experience}}
			\begin{itemize}
				\item Engineered shared, centrally maintained YAML pipeline templates replacing Classic pipelines; migrated 6 projects, enabled parallel product edition builds and standardized pipelines for new projects via documenting and enforcing rules via PR validation policies
				\item Reduced post-build cleanup from up to 15 minutes to < 10 seconds; scheduled agent cleanup to prevent disk-full failures
				\item Configured Windows Server build agents to keep interactive sessions for running UI smoke tests I introduced to build pipelines in order to catch application startup failures before manual QA
				\item Implemented centralized PR validation template via solution builds or unit tests, reporting for outdated internal dependencies, manual testing build links reporting to PR. Documented QA/release workflows and setups for consumer pipelines to support adoption of centralized YAML templates across the team
			\end{itemize}

			{\bf{Desktop}}

			\begin{itemize}
				\item Built a commercial WinUI 3 Bluetooth battery display application, analyzed low-level communications and protocols via Wireshark. Used MVVM, ReactiveUI and ensured device fetching logic is testable
				\item 6,000-line single-file screenshot editor split into a viewer and an annotation editor built on reusable tools and command history; delivered in increments so users got a more stable editor early
				\item Legacy desktop apps migrated from .NET Framework, .NET Core 3.0 to modern .NET versions and SDK-style projects
			\end{itemize}			
		}
	\entry
		{04.2023 -- 06.2024\\\footnotesize{part time}}
		{Mentor}
		{Professions of the Future Academy}
		{\texttt{C\#}\slashsep\texttt{Unity}\slashsep\texttt{Python}\slashsep\texttt{Lua}}
		{\location{Romny, Sumy region, Ukraine}}
		{}
	\entry
		{09.2023 -- 05.2024\\\footnotesize{part time}}
		{Mentor (Section Leader)}
		{Junior Academy of Sciences of Ukraine}
		{\texttt{C\#}\slashsep\texttt{C++}}
		{\location{Romny, Sumy region, Ukraine}}
		{}
\end{entrylist}

%----------------------------------------------------------------------------------------
%	EDUCATION
%----------------------------------------------------------------------------------------

\cvsect{Education}

\begin{entrylist}
	\entry
		{2024 -- present \\ \footnotesize remote}
		{Bachelor's Degree}
		{Cherkasy State Technological University}
		{Computer Science}{}{}
	\entry
		{2020 -- 2024}
		{Junior Bachelor's Degree}
		{Romny Professional College of Vadym Hetman Kyiv National Economic University}
		{Computer Engineering}{}{}
\end{entrylist}

%----------------------------------------------------------------------------------------
%	ADDITIONAL INFORMATION
%----------------------------------------------------------------------------------------

\begin{minipage}[t]{0.27\textwidth} % Left column width
	\vspace{-\baselineskip} % Required for vertically aligning minipages
	\cvsect{Languages}
	
	\textbf{Ukrainian} - native\\
	\textbf{English} - B2 (upper-intermediate)\\
\end{minipage}
\hfill % Automatic horizontal whitespace
\begin{minipage}[t]{0.7\textwidth} % Right column width
	\vspace{-\baselineskip} % Required for vertically aligning minipages
	\cvsect{Skills \& Technologies}

\textbf{Languages:} C\# $\bullet$ Python $\bullet$ C++ $\bullet$ TypeScript

\textbf{Web:} REST APIs $\bullet$ ASP.NET Core $\bullet$ EF Core $\bullet$ PostgreSQL $\bullet$ FastAPI $\bullet$ SignalR $\bullet$ Blazor

\textbf{Frontend \& UI:} WinUI 3 $\bullet$ WPF $\bullet$ UWP $\bullet$ ReactiveUI $\bullet$ Rx.NET $\bullet$ Angular $\bullet$ React

\textbf{Tools:} Azure DevOps $\bullet$ CI/CD $\bullet$ Docker $\bullet$ YAML $\bullet$ MSBuild $\bullet$ PowerShell $\bullet$ Bash
\end{minipage}

%----------------------------------------------------------------------------------------

\end{document}
